\section{Global geometry and finite-sample selection}
\label{sec:theory}

We first establish that the population mixing risk can have multiple valleys, then certify its global candidates, and finally study how finite samples shift or select the optima. Throughout, $w$ denotes the mixture weight $\lambda$ of Sec.~\ref{sec:setup}, and $R_n(w)$ is the actual full-ridge target risk. Proofs and exact certificates are in the supplementary material.

\subsection{Two coordinates suffice for multiple valleys}
\label{sec:population-geometry}

Taking the population limit of the full-ridge error in Eq.~\eqref{eq:full-error} gives the following risk; the derivation is provided in supplementary Theorems T1-b and T1-c. With $D_j(w)=(1-w)a_j+wb_j+\mu$,

\begin{equation}
 R_0(w)=\sum_j a_j\frac{(wb_j\delta_j-\mu\beta_j^*)^2}{D_j(w)^2}.
 \label{eq:r0-coordinates}
\end{equation}

\begin{proposition}[Minimum dimension for multiple valleys]\label{prop:multimodal}
For the population risk in Eq.~\eqref{eq:r0-coordinates}, when $p=1$, $R_0'(w)=L(w)/D(w)^3$ with affine $L$, so there is at most one nondegenerate local valley. In contrast, a certified two-dimensional example satisfying Assumptions~\ref{ass:a1}--\ref{ass:a3} has three simple interior stationary points
\begin{equation}
 w_1\approx0.264842,\quad w_2\approx0.727011,\quad w_3\approx0.869112,
\end{equation}
which are respectively a local minimum, a local maximum, and a local minimum. Interval risk comparison proves that $w_1$ is the unique global minimizer. Thus $p=2$ is the minimum dimension for multiple valleys.
\end{proposition}

Figure~\ref{fig:teaser} (left) shows this certified double-valley risk curve. Varying its second-coordinate weight $t$ gives valley-pair creation at $t_1\approx0.9375848$ and annihilation at $t_2\approx1.2003419$. Figure~\ref{fig:teaser} (right) tracks the risks at the stationary points as $t$ varies. Between these critical values, the global optimum switches branches, and at some $t_*\in[1.06085925,1.06085926]$ the two local minima have equal risk: two different weights are global minimizers.

\subsection{Certifying global population optima}
\label{sec:global-certification}

Multiple valleys require a comparison of all candidate optima. Let $L_j(w)$ be the affine numerator of the derivative of $(wb_j\delta_j-\mu\beta_j^*)^2/D_j(w)^2$. Then $R_0'(w)=\sum_j a_jL_j(w)/D_j(w)^3$. Clearing the positive denominators gives

\begin{equation}
 P(w)=\sum_j a_jL_j(w)\prod_{k\ne j}D_k(w)^3=0,
 \qquad \deg P\le3p-2.
 \label{eq:stationary-polynomial}
\end{equation}

For fixed parameters, isolating all real roots of $P$ in $(0,1)$ and comparing them with $w=0,1$ certifies the global minimizers without assuming convexity or unimodality. Degenerate roots and flat cases are handled separately.

Zero mismatch gives a tractable special case in which the shape itself determines how to find the optimum.

\begin{proposition}[Log-convexity and flatness at zero mismatch]\label{prop:delta-zero-shape}
Suppose $\delta=0$ and let $\mathcal J=\{j:a_j\beta_j^{*2}>0\}$. Then $R_0$ is log-convex on $[0,1]$. If some $j\in\mathcal J$ has $a_j\ne b_j$, then $R_0$ is strictly log-convex and has a unique minimizer. If every $j\in\mathcal J$ has $a_j=b_j$, or $\mathcal J$ is empty, then $R_0$ is constant and every weight is optimal. In the strictly log-convex case, the signs of $R_0'(0^+)$ and $R_0'(1^-)$ determine whether the unique optimum is at zero, one, or an interior point.
\end{proposition}

Thus the same model includes both a certified multivalley example and a zero-mismatch regime with a unique optimum or a flat risk. These population cases provide distinct starting points for finite-sample analysis.

\subsection{From population geometry to finite samples}
\label{sec:finite-selection}

To relate the certified population candidates to actual-risk optima, we need uniform control of the risk, its slope, and its curvature.

\begin{theorem}[Second-order full-ridge risk, T1-c]\label{thm:t1c}
Consider the standard full-ridge estimator under the Gaussian design and noise model, with fixed $p$, fixed $\mu>0$, and fixed $\rho=\nh/\nr$. Then, uniformly for $w\in[0,1]$,
\begin{equation}
  R_n(w)=R_0(w)+\frac{R_1(w)}{\nr}+o(\nr^{-1}),
  \label{eq:t1c-main}
\end{equation}
where $R_0$ and $R_1$ are explicit matrix coefficients derived in the supplementary. The first two derivatives satisfy the same uniform expansion.
\end{theorem}

The expansion applies to actual full ridge at fixed dimension. Its uniform derivatives preserve positive local curvature for sufficiently large samples and permit the following three conditional selection results.

\begin{theorem}[Nondegenerate interior optimum]\label{thm:optfull}
Suppose $R_0$ has a unique global minimizer $w_0\in(0,1)$ with $R_0''(w_0)>0$. Then, for all sufficiently large $\nr$, $R_n$ has a unique global minimizer $w_n^*$ on $[0,1]$, with
\begin{align}
 w_n^*&=w_0-\frac{R_1'(w_0)}{\nr R_0''(w_0)}+o(\nr^{-1}),\notag\\
 R_n(w_n^*)&=R_0(w_0)+R_1(w_0)/\nr+o(\nr^{-1}).
 \label{eq:optfull-main}
\end{align}
Consequently, $\mathcal G_n=R_n(0)-R_n(w_n^*)$ satisfies $\mathcal G_n=\mathcal G_0+\mathcal G_1/\nr+o(\nr^{-1})$, where $\mathcal G_0=R_0(0)-R_0(w_0)$ and $\mathcal G_1=R_1(0)-R_1(w_0)$.
\end{theorem}

\begin{theorem}[$R_1$ selects among multiple population optima]\label{thm:branchselect}
Suppose $R_0$ has exactly two global minimizers $0<w_1<w_2<1$, both with positive local curvature. If $R_1(w_1)<R_1(w_2)$, then for all sufficiently large $\nr$ the global optimum is attained only near $w_1$, and $w_n^*\to w_1$. The risk gap between the two local branches is $(R_1(w_2)-R_1(w_1))/\nr+o(\nr^{-1})$. If the first-order values tie, this expansion does not select a branch.
\end{theorem}

\begin{theorem}[Flat $R_0$ at matched spectra and zero mismatch]\label{thm:flatR1}
Let $\Sigma_r=\Sigma_h=\Sigma$ and $\delta=0$. Then $R_0$ is constant and there are explicit nonnegative coefficients $A,B$ such that
\begin{equation}
 R_1(w)=A(1-w)^2+Bw^2.
 \label{eq:flat-r1}
\end{equation}
If $A+B>0$, then $R_n$ is strictly convex for sufficiently large $\nr$, $w_n^*\to w_1^*=A/(A+B)$, and
\begin{align}
 R_n(w_n^*)&=C+\frac{AB}{A+B}\frac1\nr+o(\nr^{-1}),\notag\\
 \mathcal G_n&=\frac{A^2}{A+B}\frac1\nr+o(\nr^{-1}),
 \label{eq:flat-r1-optimum}
\end{align}
where $C=R_0$ and $\mathcal G_n=R_n(0)-R_n(w_n^*)$. If $A=B=0$, then $\beta^*=0$ and both label-noise variances vanish, so $R_n(w)=0$ for every sample size and weight: all weights are optimal and $\mathcal G_n=0$ exactly.
\end{theorem}

\paragraph{Implications for mixing.}
A positive zero-mixing derivative cannot by itself justify discarding the source, and a unique global minimizer does not imply unimodality. Global population certification identifies the candidates; the uniform expansion then determines local shifts and first-order selection under the conditions above. It does not classify the entire finite-sample landscape at an arbitrary sample size.

\paragraph{Further consequences.}
The supplementary material derives matched-spectrum gain formulas, a first-order exchange approximation, and a dosage rule for an explicitly defined quadratic surrogate. The dosage rule concerns that surrogate; the exchange approximation is restricted to first-order risks and a fixed-regularization baseline.
